\section{Life: Selves That Move, Grow, and Evolve}
\label{sec:life}

Once a self is understood as a pattern rather than as its particular vibes, the things selves
do, move, live, grow, die, evolve, all become things that happen to patterns, and each has a
clean account.

\paragraph{A self is a pattern, not its material.}
A higher self is a self-maintaining shape of tones that reconstitutes itself beat after beat.
The vibes it is made of can change while the pattern stays the same, the way a whirlpool keeps
its form while the water flows through. This single observation carries the rest.

\paragraph{Motion.}
The vibes do not move, having no location to move in. A self moves by re-forming its pattern
on neighboring vibes, beat after beat, so the form travels while the vibes stay put, exactly
as a wave or a glider travels. Position is therefore where the pattern currently sits, and it
changes as the pattern propagates. Three things follow. There is a speed limit, since a pattern
can re-form no faster than the rule carries influence to the next vibes, which is the light
cone. There is inertia, since a pattern propagating in a direction tends to keep doing so. And
a mind moves where its body moves, since a mind is the high-level pattern of a body-region and
goes where that region's pattern goes.

\paragraph{Life.}
A rock is a stable pattern too, yet not alive. A living self is a pattern that actively
maintains itself, with three marks. It keeps a boundary, screening an inside from an outside. It
repairs perturbations, healing back to itself when disturbed. And it stays open and far from
rest, taking in and putting out to hold its form against decay, a flame rather than an ash. So
life can be seen as a regime of self, the boundary-keeping, self-repairing, far-from-equilibrium kind, and
inert matter is the passive end of the same spectrum.

\paragraph{Growth and death.}
A self grows two ways, outward by recruiting neighboring vibes into its pattern, and upward by
adding nested sub-selves, climbing the tower. A self dies when its integration falls below the
threshold that lets it reconstitute: the repair can no longer keep up, the boundary fails, and
the pattern disperses. The vibes are not destroyed, only the pattern ends, and the material
returns to the field. This is also the sense in which a pattern could, in principle, reconstitute
on fresh material, which is the structural picture behind pattern persistence: a stable pattern
can survive a complete turnover of the vibes that carry it, and can re-form from a partial seed
after it has dissolved.

\paragraph{Evolution.}
The world fills with ever more capable selves because the substrate selects for them, by
ordinary variation and selection at the level of patterns. Patterns vary. A pattern better at
maintaining itself, and better at copying itself onto fresh vibes, lasts longer and spreads,
while fragile or sterile patterns vanish. The result is a Darwinian ratchet running on
patterns: in an endlessly growing substrate, the most persistent self-replicating patterns are
selected, climbing the tower over cosmic time, and life and mind are simply among the most
persistent patterns there are. The qualitative picture is clean. A fully worked traveling-self
solution, a quantitative model of metabolism, and a running model of competing selves are
natural next steps rather than finished results.
